% 本节专属的素材：图放本章目录的 figures/03-board/，代码放 code/03-board/。
% 章内多节共用的放 figures/ 或 code/ 根下；跨章复用的放在主用它的那一章。
% 正文里按相对本章目录的路径引用，例如 \includegraphics{figures/03-board/xxx}。
\section{开发板与实验平台}
\label{sec:02-digital-logic-03}

% 本节是全书的实验平台说明：介绍本书统一使用的开发板，以及"写码 → 上板 → 看现象"
% 这条链路上要用到的板级知识。各章实验都默认在这块板子上完成，因此本节只写一遍，
% 其它章节需要时用 \cref{sec:02-digital-logic-03} 引用，不重复介绍板子。
%
% 资料来源：《竞业达开发板使用说明》（课件，含板卡照片、板上器件说明与引脚分配查询）。
% 落地前请核对课件版本与板卡丝印，器件数量、有效电平等以板卡实物和官方引脚表为准。
%
% TODO：撰写本节正文。建议结构：
%   1) 问题引入与学习目标：为什么先认识板子；学完能独立把一份设计下载到板上并看到现象。
%   2) 板卡整体结构与主要接口：型号 \BoardModel{}（\BoardFPGA{}），主要接口与器件分布配图。
%   3) 板上器件怎么用：
%        · 控制器件：64 个拨动开关（分 8 组；向上断电=逻辑 0，向下通电=逻辑 1）；
%          8 个按键（按下断电=0，松开通电=1）
%        · 显示器件：32 个 LED（分 4 排，高电平点亮）；
%          4 位共阴极数码管（段选 LEDx_A–LEDx_G、LEDx_DP 高有效，位选 LEDx_CS 低有效）；
%          两位数码管共用一套段码，靠位选以约 10 MHz 的频率分时选通，视觉上"同时"显示
%        · 发声器件：无源蜂鸣器，需要方波驱动，频率越高音调越尖
%      段码表（共阴、高有效，SEG[7:0] 自高位到低位为 DP,G,F,E,D,C,B,A）：
%        0→0x3F  1→0x06  2→0x5B  3→0x4F  4→0x66
%        5→0x6D  6→0x7D  7→0x07  8→0x7F  9→0x6F
%   4) 上板准备与下载：JTAG 转接板、JTAG 盒与 TYPE-C 线的连接顺序，接口孔数要先对上，
%      插拔不要用蛮力；在 Vivado 里识别设备并下载比特流。
%   5) 引脚分配查询：用器件—引脚对照表（课件中的 function pin 与 xc7k325t 两份表）查
%      XDC 约束里要写的引脚名，以一个 LED 加一个拨动开关为例走一遍。
%   6) 小结与思考题：把"仿真 → 上板"这条闭环跑通一次。
本节待撰写。
